\begin{abstract}
Every regular language over a finite alphabet has generalized star height
at most three, with complement taken in the same free monoid. We express
finite-monoid computations using a prefix code of word pieces.
\end{abstract}

\section{Introduction}\label{sec:introduction}

The star operation permits a finite language expression to describe an
arbitrarily long repetition. Its nesting depth measures how many layers
of repetition an expression uses. Once complement is allowed, Boolean
conditions can replace some of those layers. The question studied here
is whether a fixed number of layers suffices for every regular language,
independently of its alphabet and recognizing automaton.

Fix a finite alphabet $\Sigma$. A \emph{generalized regular expression}
over $\Sigma$ uses $0$, $1$, the singleton letters of $\Sigma$, union,
concatenation, complement, and Kleene star. The constants denote
$\emptyset$ and $\{\eps\}$, and complement is always taken in $\Sigma^*$.
Its height is determined by
\[
\begin{aligned}
 \sht(0)=\sht(1)=\sht(a)&=0,\\
 \sht(P\cup Q)=\sht(PQ)&=\max\{\sht(P),\sht(Q)\},\\
 \sht(\neg P)&=\sht(P),&
 \sht(P^*)&=1+\sht(P).
\end{aligned}
\]
For a regular language $L\subseteq\Sigma^*$, write $h_\Sigma(L)$ for
the least height of an expression defining $L$. In particular,
complement preserves the nesting already present in its argument. Finite
Boolean combinations preserve a common height bound, and the full
language $\Sigma^*=\neg0$ has height zero.

\begin{theorem}\label{thm:main}
For every finite alphabet $\Sigma$ and every regular language
$L\subseteq\Sigma^*$,
\[
                            h_\Sigma(L)\le3.
\]
\end{theorem}

The expressions in the theorem use the letters of $\Sigma$ itself. Their
finite parameters and their lengths may depend on the language. The
construction gives an upper bound on nesting, without a bound on
expression size or running time.

